\documentclass[12pt]{article}
\usepackage[utf8]{inputenc}
\usepackage[top=1in, bottom=1in, left=1in, right=1in]{geometry}
\usepackage{times} % Use Times New Roman font
\usepackage{setspace} % For line spacing
\usepackage{titlesec} % For section formatting
\usepackage{graphicx} % For including graphics
\usepackage{fancyhdr} % For headers and footers
\usepackage{hyperref} % For hyperlinks

% Adjust headheight and topmargin to resolve the fancyhdr warning
\setlength{\headheight}{14.5pt}
\addtolength{\topmargin}{-2.5pt}

% Line spacing
\setstretch{1.5}

% Title and section formatting
\titleformat{\section}{\normalfont\bfseries\uppercase}{\thesection}{1em}{}
\titleformat{\subsection}{\normalfont\bfseries}{\thesubsection}{1em}{}
\titleformat{\subsubsection}{\normalfont\itshape}{\thesubsubsection}{1em}{}

% Header and Footer
\pagestyle{fancy}
\fancyhf{}
\fancyhead[L]{\leftmark}
\fancyhead[R]{\thepage}

% Title Page
\title{NIH SBIR Phase I Grant Proposal}
\author{Your Name \\ Your business \\ Email: your.email@example.com}
\date{\today}

\begin{document}

\maketitle

\newpage

\section*{Abstract}
This section provides a concise summary of the proposed project. It should include the project's objectives, methods, expected outcomes, and potential impact. The abstract is often limited to one page.

\newpage

\section{Specific Aims}
State the specific goals of the project, typically in bullet points. Be clear and concise, focusing on what you intend to accomplish.

\newpage

\section{Research Strategy}
\subsection{Significance}
Explain the importance of the problem or critical barrier to progress in the field that the proposed project addresses.

\subsection{Innovation}
Describe any novel theoretical concepts, approaches, methodologies, instrumentation, or interventions to be developed or used.

\subsection{Approach}
Discuss the overall strategy, methodology, and analyses to be used to accomplish the specific aims. Include how the data will be collected, analyzed, and interpreted.

\newpage

\section{Budget Justification}
Provide a detailed justification for each budget category, explaining the necessity and reasonableness of the costs.

\newpage

\section{Biographical Sketches}
Include biographical sketches of key personnel, following the NIH format. Each biosketch is typically limited to five pages.

\newpage

\section{Facilities and Other Resources}
Describe the facilities and resources available for conducting the proposed research. This section should demonstrate that you have access to the necessary infrastructure.

\newpage

\section{References}
Enter your references manually or use refer to a .bib file in your project.

% Manually input references here. For example:
\begin{itemize}
    \item Author1, A. A., Author2, B. B. (Year). Title of the article. \textit{Journal Name}, Volume(Issue), pages.
    \item Author3, C. C., Author4, D. D. (Year). Title of the book. \textit{Publisher}.
    % Add more references as needed.
\end{itemize}

% Use a bib for references here:
%\bibliographystyle{plain}
%\bibliography{references} 

\end{document}

